\scanpage{8}
\mspage{4)}
Cette composition est associative à iso can. près, et elle a des identités bilatères (\deleted{\unsure{par}} un iso can. !) savoir
\diagram{\begin{tikzcd}
x \arrow[dr,"1"'] & x \arrow[d,"1","\sim"'] \\
& x
\end{tikzcd}} noté $\uline{1}_x$.

On a d'autre part des \deleted{\unsure{appl.}} foncteurs
\[\operatorname{Hom}_C(x,y)_{\text{discret}}\longrightarrow\uline{\mathrm{Ch}}_0(x,y)\]
\deleted{\unsure{compatibles à}}
\[f\longmapsto(f,\operatorname{Id}_{\unsure{x|y}})\]
compatibles aux compositions. Ceci posé soit
\[\pi_0(x,y)\overset{\mathrm{def}}=\pi_0\uline{\mathrm{Ch}}_0(x,y);\]
alors la composition des $\uline{\mathrm{Ch}}_0$ définit une composition des $\pi_0(x,y)$, qui est associative (str. !) et admet les classes des $\uline{1}_x$ comme unités bilatères. D'où une catégorie, que je note $\widetilde C^{\unsure{2|?}}$, et un foncteur
\[C\longrightarrow\widetilde C.\]

Pour l'instant, on n'a utilisé rien sur $(C,W)$, [\deleted{\unsure{ni}} \unsure{ni} la donnée des $\mathrm{cof}\subset\mathrm{Fl}\,C$; \unsure{en se contentant} de prendre les flèches $i$ dans (\deleted{$\mathrm{cof}\cap W$}\inserted{$\unsure{\mathrm{Couniv.}W}$}) --- mais \deleted{\unsure{si}} on peut garder la construction, avec $\mathrm{cof}$ \inserted{\unsure{pas néc.}\ $=\mathrm{cof}_W$,} tel que $\mathrm{cof}\cap W\subset\deleted{\unsure{W}}\unsure{\mathrm{Couniv.}W}$]. \deleted{On} \inserted{Mais,} on veut prouver que ce foncteur transforme $f\in W$ en iso.
\diagram{\begin{tikzcd}
x \arrow[r,"f","\sim"'] & y \arrow[dr,"g?"'] & x \arrow[d,"i?","\sim"'] \\
&&\bar x
\end{tikzcd}}
on cherche $g,i$ tels que $(gf,i)\in\text{comp. conn. de }\uline{1}_x$ dans $\pi_0(\uline{\mathrm{Ch}}_0(x,x))$ et $(\bar f g,j)\in\text{comp. conn. de }\uline{1}_y$.
\diagram{\begin{tikzcd}
& x \arrow[r,"f","\sim"'] \arrow[d,"i","\sim"'] & y \arrow[d,"j"',"\sim"] \\
y \arrow[r,"g"'] & \bar x \arrow[r,"{\bar f\quad\sim}"',"\sim"] \arrow[ur,phantom,"(\text{cocart})" description] & \bar y
\end{tikzcd}}
\mspage{7}
